\documentclass[11pt]{article}
\usepackage[utf8]{inputenc}
\usepackage[margin=1.15in]{geometry}
\usepackage{amsmath,amssymb}
\usepackage{array}
\usepackage{longtable}
\usepackage{setspace}
\usepackage{hyperref}
\usepackage{titlesec}
\setstretch{1.1}
\hypersetup{colorlinks=true,linkcolor=black,citecolor=black,urlcolor=black}
\titleformat{\section}{\large\bfseries}{\thesection.}{0.6em}{}
\titleformat{\subsection}{\normalsize\bfseries\itshape}{\thesubsection.}{0.6em}{}
\newenvironment{principle}{\begin{quote}\itshape}{\end{quote}}

\title{Histories That Rewrite Their Own Possibilities\\[0.4em]
\large Organization Without a Prior Plan}
\author{Flyxion\\Independent Researcher}
\date{September 2026}

\begin{document}
\maketitle

\begin{abstract}
The finished structure of a system need not exist at the beginning, even as information. Each stage need only modify the admissible continuations of the next. This essay defends that claim across a chain of cases in which organized outcomes invite the inference of an organizing plan: the repeated lone messenger of Job 1, the interrupted transmission at Mark 16:8, gravitational structure formation, mineral evolution, autocatalytic chemistry, biological niche construction, cultural inheritance, and Nietzschean genealogy. The result is stronger than the familiar claim that complexity emerges from randomness. Variation, constraint, repeated trials, differential persistence, feedback, and time can yield organization, and when outcomes alter the conditions that govern later outcomes, the possibility space itself changes as the process runs. Claims are graded throughout as established findings, plausible mechanisms, illustrative analogies, or the essay's own synthesis, and an appendix table collects the grades and sources.
\end{abstract}

\section{The problem of the unplanned pattern}

Some patterns look designed because they solve a problem. A messenger arrives in time. A disk of stars rotates in a plane. A reaction network sustains itself. A moral code appears to follow from a principle. In each case the solved problem tempts the observer to look backward for a solver.

The temptation is not irrational. Sometimes a solver exists. Authors do plan, engineers do design, and institutions do legislate. The difficulty is that the presence of a solved problem does not by itself discriminate between a solver and a process that retained solutions while discarding failures. The method of this essay is to ask, of each case, one portable question: \emph{which stage of the process altered the conditions of the next, and did the account of the outcome require a plan that the process itself does not supply?}

\subsection{Terms}

Three words are easily run together and should be kept apart. \emph{Order} means regularity or symmetry, as in a crystal. \emph{Complexity} means a combination of differentiation and integration, as in an ecosystem. \emph{Organization} means persistent differentiation into mutually constrained components or relations, as in a galaxy with distinct components, a cell, or a narrative that delivers its information. Some organized systems, cells and narratives among them, are also functionally organized, meaning that their parts contribute to the persistence of the whole or to a purpose, but function is not required for the category. The essay's subject is organization in the neutral sense. Order and complexity appear only where they bear on it.

\subsection{Roadmap and grading}

Sections 2 to 4 treat narrative, because the constraints there are easiest to see at small scale: Job 1, Mark 16:8, and the general case of narrative without a single designer, together with the separation of literary construction from historical existence. Section 5 turns to gravity, where the law is known and the organization is still not written into the law. Sections 6 and 7 treat mineral evolution and autocatalysis, where feedback first appears. Section 8 treats biological and cultural iteration, Section 9 treats genealogy as a method, and Section 10 states the synthesis. Section 11 states what follows, Section 12 lists limits, and an appendix grades every principal claim.

Claims carry one of four grades. \emph{Established} means supported by primary texts, direct observation, or standard results. \emph{Plausible mechanism} means a proposal with real support that remains contested or unproven. \emph{Illustrative analogy} means a comparison offered to clarify structure and not as evidence. \emph{Synthesis} means the essay's own proposal. Biblical references follow standard chapter and verse.

\section{Job: the survivor as a constraint on the generator}

Job 1 reports four disasters in sequence. Raiders take the oxen and donkeys and kill the servants, fire falls from heaven and consumes the sheep and the shepherds, another raiding party seizes the camels and kills the servants, and a great wind collapses the house upon Job's children. Each report closes with the same formula, in which the messenger says that he alone has escaped to tell (Job 1:15, 16, 17, 19). This much is \emph{established} by the text.

Read as four independent events, the configuration is strained. Each catastrophe is nearly total, yet each leaves exactly one survivor, and each survivor reaches Job at the right moment. If each event were a draw from the space of ordinary disasters, the recurrence of this configuration four times would be improbable.

Read as narrative construction, the improbability is a signature and not a puzzle. The story requires that Job learn of each loss, so each catastrophe must be almost total while preserving one transmission channel. Within the fictional world the survivor is of course an outcome of the disaster, since someone happened to escape. At the level of the text, however, his survival is constrained by the narrative's need for transmission: the structure that will carry the report places a requirement on the disaster. The apparent improbability therefore exposes a generative constraint. Whatever process produced this text was subject to a condition that ordinary events are not subject to. This reading is an \emph{illustrative analogy} for the general problem and a modest literary claim about the text.

An observed regularity, here the lone survivor, reveals the constraint under which the sequence was generated. It does not by itself reveal whether a single author imposed the constraint or whether a tradition of retelling converged on it.

\section{Mark 16:8: transmission narrated and transmission left open}

The earliest complete Greek manuscripts of Mark, Codex Sinaiticus and Codex Vaticanus, end at 16:8, and other early witnesses agree. Longer and shorter endings appear in later manuscripts, and the textual history is discussed in Metzger's commentary and in the continuing dispute over the longer ending \cite{metzger,head,snapp}. That the earliest attested text ends at 16:8 is \emph{established}. Whether Mark intended to end there is not, and the essay takes no position on it beyond noting the consequences of the ending as it stands.

At that ending the women flee the empty tomb, trembling and amazed, and say nothing to anyone, for they were afraid. Taken strictly, the text says only that they said nothing at that point, and a later telling remains possible. The contrast with Job should therefore be stated with care. Job narrates each act of transmission. Mark leaves it unnarrated. Job repeatedly runs
\begin{equation*}
\text{catastrophe} \rightarrow \text{survivor} \rightarrow \text{escape} \rightarrow \text{testimony},
\end{equation*}
and Mark ends with
\begin{equation*}
\text{discovery} \rightarrow \text{witnesses} \rightarrow \text{escape} \rightarrow \text{silence}.
\end{equation*}
In Job, escape is the mechanism that guarantees transmission. In Mark, escape is followed by silence, and the reader is left to supply the propagation. Read structurally, the function of escape is inverted, from guarantor of transmission to its interruption, and the narrative does not collapse. This is a literary reading of the text as it stands and not a claim about historical events or about authorial intent.

The later Gospels supply it. In Matthew the women run to tell the disciples and Jesus meets them (Matthew 28:8--10), and the promised meeting in Galilee is narrated (Matthew 28:16--20). In Luke the women report to the eleven and the rest, who take the report for nonsense (Luke 24:9--11). The longer ending of Mark supplies appearances and proclamation. That later texts supply what Mark leaves unnarrated is \emph{established}. That they did so in order to repair a felt gap is a \emph{plausible mechanism}, since the motives of the writers must be inferred.

\subsection{Identity as an established relation}

The ending belongs to a wider question about how identity is established in the Markan Passion narrative, as opposed to assumed. Judas's kiss functions as an explicit identification procedure, a signal agreed in advance so that the arresting party can pick out its target (Mark 14:44--45). A young man, a \emph{neaniskos}, follows Jesus wearing a linen cloth over his naked body, is seized, and escapes by leaving the cloth in his captors' hands (Mark 14:51--52). A young man, again a \emph{neaniskos}, sits in the tomb in a white robe and instructs the women to tell the disciples, and Peter, that Jesus goes before them into Galilee (Mark 16:5--7). Some interpreters connect the two figures, for example Scroggs and Groff \cite{scroggs}, and the essay does not depend on any identification. What can be said without one is that the second young man is the only figure in the scene who issues a message for transmission, and that the transmission is not narrated as completed.

The garment escape has a clear structural relative in Genesis 39:12, where Joseph flees from Potiphar's wife and leaves his garment in her hand. The pattern is
\begin{equation*}
\text{seizure} \rightarrow \text{retained garment} \rightarrow \text{escaped person}.
\end{equation*}
The pursuer captures the identifying exterior and fails to capture the person. The textual parallel is \emph{established}, and the inference of literary dependence is an \emph{illustrative analogy}. A fleeing figure whose clothing is seized and left behind is also a familiar farcical device. Keaton and St.\ Clair's \emph{The Goat} (1921) is a later comic treatment of a related problem, since its plot turns on identity fixed by a token, a manipulated photograph on wanted posters, that attaches the wrong man to a crime \cite{goat}. Judas's kiss addresses the same design problem from the other side. In each case identity is carried by an exterior mark that can be separated from the person. Performance matters here because darkness, costume, entrances, exits, and exchanged garments make identity genuinely ambiguous to an audience in a way that prose can suppress.

Qur'an 4:157, with its \emph{shubbiha lahum} tradition, shows at minimum that uncertainty about what happened to Jesus at the crucifixion became a live interpretive option in the history of Muslim thought \cite{lawson}. Costume, darkness, or misidentification are sufficient to generate such ambiguity without any supernatural change of appearance. This is a reconstruction of narrative mechanics and not evidence of what historically occurred.

These observations support a purely literary question that requires no historical conclusion. At what points does the text establish that the person at a later time $t_2$ is the same person as at an earlier time $t_1$, and at what points does the reader silently supply that identity relation? A narrative is partly a device for managing exactly this bookkeeping, and its gaps show where the bookkeeping is left to the audience.

\section{Narrative without a single designer}

A narrative need not have a single designer any more than a physical structure does. Authors inherit motifs, and performers discover devices that work on audiences. Communities preserve some versions and lose others. Contradictions provoke repairs, and later authors regularize ambiguities left by earlier ones. Institutions stabilize particular variants and exclude others. The eventual cultural object can possess far more organization than any participant planned. Oral-formulaic studies document the reuse of inherited formulas and themes in traditional performance \cite{lord}, and the general framework of cultural evolution treats transmission, variation, and differential retention as processes with population-level consequences \cite{boyd,mesoudi}. These sources support the \emph{plausibility} of the mechanisms. They do not show that any given text arose by them.

The alternatives are therefore not exhausted by ``this happened as narrated'' and ``one author consciously invented it.'' One possible mechanism for the lone messenger of Job is differential retention, in which versions that deliver their information are copied and versions that do not are lost. The essay supplies no textual or traditional evidence that this occurred, and the constraint of Job 1 is equally compatible with a single author who saw the requirement directly. For Mark 16:8 the record differs in kind, because the later endings and the parallel Gospel accounts are themselves evidence that the gap was addressed by later hands, although the motives of those writers remain a matter of inference. In both cases the organization of the finished text is real, and retention, repair, and relocation are candidate sources of it that need no representation of the outcome in advance.

This is a claim about what the generative process must contain. It does not deny that individual authors made deliberate choices, and many did. It denies only that the global structure of a tradition must be traced to a single intention.

\subsection{Literary construction and historical existence}

The evidential point that follows concerns two variables that are easily conflated. Richard Carrier's historicity argument applies Bayesian reasoning to the question of what reference class should fix the prior probability assigned to the Gospel Jesus, comparing him with mythologized heroic figures \cite{carrier1,carrier2}. Its reference-class construction and probability assignments are disputed in detail \cite{gullotta}, and nothing here depends on its numerical conclusions. What transfers is the reference-class question, together with a separation that the essay offers as its own \emph{synthesis}.

Joshua, Moses, Samson, Elijah and Elisha, the Psalms, Isaiah, and other inherited texts form a store of narrative machinery that can be relocated, morally reinterpreted, and adapted to new audiences. That a later protagonist resembles earlier literary figures does not establish nonhistoricity, since historical people can be mythologized. Nor does it establish historicity. The resemblance is evidence about how the narrative was built and is silent, taken alone, about whether an underlying person existed. Recognizing biblical or classical heroic structures in Superman would tell us about the construction of Superman and would supply no independent attestation of Clark Kent. Finding Heraclean motifs in narratives about Alexander would not erase the independent evidence for Alexander. Literary construction and historical existence are separate variables, and the productive question for any episode is which evidence is independent of the literary system under examination.

The method that follows is more discriminating than cataloguing parallels. For each episode one asks which earlier narrative operations it reproduces, what has been transformed, which details are required by the new story and not merely available, how identity and testimony are established, and what historical residue remains once material explicable by literary construction has been subtracted. The residue may be large, small, or empty.

\section{Gravity: structure without a representation of structure}

The narrative cases show that organization can plausibly accumulate through retention and repair without a planning agent. The question now is whether the same structure appears where the generative process is fully known and fully impersonal. Gravitational dynamics supply the cleanest case, because the governing law is simple, the initial conditions can be specified exactly, and nothing in the law refers to the outcome.

Gravity acting on a nearly uniform distribution of matter amplifies small density perturbations. Regions that are slightly denser attract more, become denser still, and draw material from their surroundings, while neighboring regions are depleted. Cosmological simulations and theory describe the result as a web of filaments, sheets, and nodes surrounding underdense voids \cite{bond}. Concentration and depletion arise together and are complementary consequences of redistribution. A void is not a place where nothing happened. Its relative emptiness is partly the consequence of matter having accumulated elsewhere. These points are \emph{established} in cosmology, and the void review cited below notes that voids are underdense regions that do contain galaxies \cite{vdw}.

In simple gravitational experiments with a highly symmetric initial distribution, matter can collapse coherently and rebound, behaving almost like a single ball, and the resulting organization is trivial because the initial state was already nearly uniform. Introduction of heterogeneity produces differentiation at many scales. This is offered as an \emph{illustrative analogy} drawn from simple simulation, without citation, and the argument does not rest on it.

\subsection{Necessary qualifications}

Gravity does not generate net angular momentum from a perfectly symmetric state of zero angular momentum. In the standard account, galaxies acquire angular momentum through tidal torques exerted by surrounding matter on growing perturbations, which requires asymmetry in the initial distribution \cite{peebles,white}. Galaxy morphology further depends on velocity dispersion, gas dissipation, mergers, and other processes. None of this weakens the conceptual point. The claim is not that gravity alone produces galaxies. It is that a simple local law applied to heterogeneous initial conditions can generate macroscopic organization without that organization being represented beforehand.

This yields the first general principle, graded \emph{established} for gravity and \emph{synthesis} as a general statement:
\begin{principle}
Structure does not entail a prior representation of the structure.
\end{principle}

\section{Mineral evolution: a changing chemical space}

Hazen and colleagues propose that the mineral inventory of a planet is not fixed but evolves. Their framework describes ten stages, beginning with roughly a dozen minerals in the pre-planetary nebula and proceeding through planetary differentiation, plate tectonics, and the influence of life on atmosphere and oceans to the more than five thousand mineral species now recognized, and they attribute a large share of present diversity to biological effects on surface environments \cite{hazen2008}. Whatever the detailed staging, the central observation is \emph{established} as a framework: the set of minerals present at one epoch is a consequence of earlier history, and it in turn conditions later chemistry. Planetary history modifies the chemical possibility space available to subsequent planetary history.

The same lineage was traced earlier, with a different aim, in \emph{From Minerals to Minds} \cite{m2m}. That essay follows differential persistence from mineral evolution and autocatalytic chemistry through membranes, swarm intelligence, and human institutions, and formalizes it as the retention of local constraint modifications that raise entropy production. Its chemical section relies on ordinary tidal-pool evaporation and drying cycles and cites Hazen's Hadean paleomineralogy \cite{hazen2013}, and it does not depend on the faster-rotation tidal proposal discussed below. The present essay covers the same chemical ground but asks a different question. The earlier claim concerns what is retained and at what rate. The present claim concerns how retained outcomes alter the space of later outcomes, and it tests that claim against narrative and genealogical cases as well as physical ones.

Mineral surfaces bear on origin-of-life scenarios. Laboratory work shows that surfaces of common oxides, silicates, and carbonates select and concentrate specific amino acids, sugars, and other molecules, and that clay surfaces can scaffold the formation of RNA oligomers of tens of nucleotides. The same review records serious caveats: growing polymers can bind ever more tightly to the surface, and most experiments use single minerals and simple solutions that do not reproduce geochemical complexity \cite{hazen2010}. Adsorption and concentration by minerals are \emph{established} in the laboratory. Their sufficiency for the origin of life is a \emph{plausible mechanism}.

\subsection{A distributed reactor}

The following picture is an \emph{illustrative analogy}, and its limits come first. It is not a claim that the Earth was a cell. Cells have stronger boundaries, regulated transport, and internally maintained organization. What the picture captures is narrower: a distributed reactor consisting of many partially isolated microenvironments, in which cracks and pores supply locality, mineral surfaces supply adsorption and catalytic selectivity, water movement supplies transport, and wet-dry cycles repeatedly concentrate and disperse material. Geological boundaries supply partial compartmentalization long before lipid membranes exist. A membraneless planetary cell is a vivid name for this arrangement and should be held loosely.

\subsection{A contested proposal about tides}

One proposal concerns tidal forcing. Lathe proposed that a closer early Moon and faster terrestrial rotation subjected coastal areas about 3.9 billion years ago to rapid flooding and drying, with a periodicity of roughly two to six hours, and that the resulting cycles of dilution and concentration, resembling the polymerase chain reaction, could favor the replication of nucleic-acid-like polymers \cite{lathe}. The proposal was challenged in a published comment \cite{varga} and answered by its author. It is a \emph{plausible mechanism} with contested standing, and the essay's argument does not depend on it. The further suggestion that abrasion at shorelines generates fresh mineral surface at a rate relevant to prebiotic chemistry is conjecture in this essay, and no source is offered for it.

\subsection{How parallelism changes the probability question}

Suppose a chemically interesting event has probability $p$ per opportunity. The next calculation assumes that the opportunities are approximately independent. Under that assumption, the probability that the event has occurred at least once after $N$ opportunities is
\begin{equation}\label{eq:parallel}
P(\geq 1) = 1-(1-p)^N \approx 1-e^{-Np},
\end{equation}
where the approximation holds for small $p$. This is standard mathematics. Applied to the mineral case, it says that a reaction that would be fantastically unlikely in a particular crack in a particular cycle need not be remotely probable there, and can nonetheless become probable somewhere when planetary geology supplies sufficiently many trials over sufficiently long periods. The value of $N$ for real prebiotic settings is not known, and the essay does not estimate it. The equation makes the negative point that a small $p$ does not by itself exclude an event, and it makes no positive claim about the origin of life.

The independence assumption fails exactly where the interesting event begins. Once a network appears that improves the conditions for its own continuation, later trials are no longer drawn from the original distribution, and Equation~\eqref{eq:parallel} no longer describes them. The calculation explains why a rare first event is not surprising. It does not describe what happens afterward, which is the subject of the next section.

\section{Autocatalysis and the first historical dependence}

Parallelism alone is not the crucial transition. The decisive change occurs when a rare reaction network modifies its environment in a way that increases the probability of its own continuation.

The formal theory of autocatalytic sets, introduced by Kauffman and developed by Hordijk and Steel, describes sets of reactions in which every reaction is catalyzed by a member of the set and every reactant can be built from a supplied food set \cite{kauffman,hordijk2004,hordijk2019}. That such sets are mathematically well defined, and that they can arise with appreciable probability in models of chemical reaction systems, is \emph{established} within the models. Whether prebiotic chemistry realized them is a \emph{plausible mechanism} under investigation. Autocatalysis in this formal sense does not require the simple textbook form $A \rightarrow 2A$, in which one molecule copies itself. It requires catalytic closure of the network as a whole.

The essay's proposal is broader than catalytic closure and should be marked as such. Standard theory requires each reaction in the set to be catalyzed by a member of the set, whereas the proposal here allows a member to raise the rate of a reaction indirectly, by changing the shared environment. Consider
\begin{equation*}
A \rightarrow B \rightarrow C \rightarrow D,
\end{equation*}
where $D$ increases the availability of $A$, accelerates the step $A \rightarrow B$, removes an inhibitor, alters local pH, modifies a mineral surface, or otherwise raises throughput around the cycle. No component needs to know that the loop exists, and no single molecule copies itself. The feedback operates through the environment and not only through direct catalysis. This environment-mediated version is the essay's \emph{synthesis}, an extension of the formal theory and not a result from it.

What must appear first, on this view, is not something resembling a modern gene but a closed positive causal path. That is a far weaker requirement, and one that a distributed set of microreactors is well placed to satisfy somewhere.

\subsection{Terminology: persistence before reproduction}

The process at this stage should be described as chemical selection, differential persistence, or selection among autocatalytic networks, and not automatically as Darwinian natural selection. Full Darwinian evolution requires reproduction with heritable variation and differential reproductive success. Networks that facilitate their own continuation persist and accumulate, and networks that do not are lost, but this is not yet reproduction in the required sense. The distinction shows that the argument does not need to wait for genes. One conceptual sequence, offered as a \emph{plausible} ordering and not as a demonstrated history, is
\begin{align*}
&\text{random chemistry} \rightarrow \text{persistent chemistry} \rightarrow \text{self-facilitating chemistry}\\
&\rightarrow \text{autocatalytic networks} \rightarrow \text{environmental modification}\\
&\rightarrow \text{compartmentalization} \rightarrow \text{heredity} \rightarrow \text{Darwinian evolution}.
\end{align*}
No stage requires a sudden import of a foreign principle called life. The interesting question becomes where increasingly powerful forms of self-maintaining continuation appear.

The important conceptual transition occurs before genes. A chemical system acquires a primitive historical dependence when its products alter the probability distribution governing its subsequent reactions. This gives the second general principle, graded \emph{synthesis}:
\begin{principle}
Once outcomes modify the conditions that produce later outcomes, history ceases to sample repeatedly from an unchanged possibility space.
\end{principle}

\section{Biological and cultural iteration}

Organisms do not merely adapt to environments. They modify them, and in doing so alter the selection pressures that act on their successors. Life has changed atmospheric chemistry, oceans, soils, and mineral diversity, and niche construction theory treats this modification as a process in evolution in its own right \cite{oslf}. The phenomena are \emph{established}. The theory's claim that niche construction is a distinct evolutionary process, and not a consequence of standard selection, is influential and debated. The essay's argument does not depend on that dispute, since the feedback structure is present even if niche construction is fully reducible to standard selection theory.

Cultural evolution supplies a further iteration, graded as an \emph{illustrative analogy}, and it connects back to the opening examples. Stories change the environment in which later stories are produced. A successful narrative structure becomes available for reuse. A contradiction can provoke harmonization, and a difficult ending can provoke continuation, as Mark 16:8 plausibly did. A theological innovation changes what subsequent writers must explain, and canonization changes which variants remain available for recombination. Each of these is an outcome that alters the conditions under which later outcomes are selected.

The analogy has a real weakness. Cultural transmission lacks a fixed replicator of the kind genes supply, and some accounts hold that cultural items are reconstructed at each transmission and shaped by attractors more than copied \cite{sperber}. The mechanisms named here, retention, repair, relocation, and stabilization, do not require replicators, but the analogy should be tested against them case by case.

\section{Genealogy as method}

Nietzschean genealogy supplies the methodological counterpart of these processes. Nietzsche repeatedly refuses to infer the origin of a value from the self-description of the finished value. A present moral concept may present itself as a product of reason, revelation, or timeless truth while having arisen through historical struggle, affective dispositions, physiological conditions, social reinforcement, and repeated reinterpretation \cite{nietzsche}.

A concrete case is his treatment of punishment in the Second Essay. Nietzsche argues that the origin of a practice and its later purpose are distinct, and that a practice that has lasted is repeatedly reinterpreted and redirected by new powers to new ends. Punishment, on his account, has no single original meaning. A cluster of meanings has accreted to a procedure that is comparatively stable, so that the procedure and its stated purpose can be traced separately (Second Essay, sections 12 and 13). The inference from a finished practice and its present justification to its origin is exactly the inference the method forbids. This is a standard reading of Nietzsche, and its application to narrative or to physical systems is an \emph{illustrative analogy}.

Genealogy therefore shares the central prohibition of the physical examples:
\begin{principle}
Do not project the finished structure backward into its origin.
\end{principle}

A galaxy does not imply a galactic blueprint. An organism does not imply that evolution anticipated the organism. An autocatalytic network does not require chemistry to have anticipated life. A mature moral system need not have originated from its present justification. And a highly organized narrative tradition does not by itself imply either a single originating author or a historical sequence matching the finished narrative. This is a prohibition on one specific inference and not a general skepticism about design. Where independent evidence shows that a plan existed, the plan is part of the explanation. The claim is only that organization is not, by itself, such evidence.

\section{The synthesis}

A first statement of the synthesis is
\begin{align*}
&\text{variation} + \text{constraint} + \text{parallelism} + \text{differential persistence}\\
&\quad + \text{feedback} + \text{time} \;\rightarrow\; \text{organized historical structure}.
\end{align*}
Each ingredient contributes a distinct function. Variation supplies candidates, constraint restricts them, parallelism supplies many trials, differential persistence retains some outcomes, feedback lets retained outcomes influence what follows, and time allows the loop to run. Table~\ref{tab:constraint} states what plays each role in the cases examined, and is an organizing summary of the \emph{synthesis} and not a set of independent findings.

\begin{table}[h]
\centering
\small
\begin{tabular}{p{0.16\textwidth}p{0.2\textwidth}p{0.27\textwidth}p{0.27\textwidth}}
\hline
Domain & Variation & Constraint & Feedback \\
\hline
Narrative & Variant tellings and retellings & Need to transmit information; audience reception; institutional canonization & Retained forms constrain later production; contradictions provoke repair \\
Gravity & Initial density perturbations & Conservation laws; local dynamics & Redistribution alters later gravitational relations \\
Mineral evolution & Chemical and physical possibilities & Thermodynamics; available surfaces and elements & New minerals and surfaces create new reaction environments \\
Autocatalysis & Molecular encounters in many microenvironments & Reaction kinetics; compartment boundaries & The network alters its own precursor availability or conditions \\
Biology & Heritable variation & Selection pressures & Niche construction alters subsequent selection \\
Culture & Variant stories and practices & Audience retention; institutional stabilization & Successful forms constrain what later authors must address \\
\hline
\end{tabular}
\caption{Roles played in each domain by the ingredients of the synthesis.}
\label{tab:constraint}
\end{table}

\subsection{Fixed and endogenous transitions}

An objection holds that this reduces to randomness plus selection, which would make it Darwinian evolution applied everywhere. Ordinary evolutionary theory is not committed to a fixed selective environment. Frequency-dependent selection, coevolution, eco-evolutionary feedback, and niche construction all make the selective regime partly a product of the population's own history. The relevant contrast is therefore between two kinds of model and not between evolution and something else. In a fixed-transition model, the rule that carries the state forward does not change:
\begin{equation}\label{eq:fixed}
S_{t+1} = T(S_t).
\end{equation}
In an endogenous-transition model, the rule is itself updated by what the process has produced:
\begin{equation}\label{eq:recursion}
S_{t+1} = T_t(S_t), \qquad T_{t+1} = F(T_t, S_{t+1}).
\end{equation}
Equation~\eqref{eq:recursion} says that the transformation regime is partly a product of previous transformations, and not merely that the state changes. Standard evolutionary theory contains endogenous-transition models. The essay's point is that this recursion is the common feature of the cases examined and that it changes what an explanation must contain. A sampling process governed by Equation~\eqref{eq:fixed}, however large the sample, cannot construct its own distribution.
\begin{principle}
History changes its own space of admissible continuations.
\end{principle}

The examples fit this schema one after another. Gravity redistributes matter and thereby changes subsequent gravitational relations. Geological evolution creates minerals and surfaces that enable chemistry previously unavailable. Autocatalytic chemistry modifies the conditions for subsequent chemistry. Biological evolution constructs niches that alter subsequent selection. Cultural evolution produces inherited forms that constrain subsequent cultural production. Genealogy reconstructs these paths without assuming that the endpoint existed implicitly at the beginning.

The essay's claim is therefore narrower than the slogans it might be mistaken for. It does not say that everything evolves, and it does not say that order comes from chaos. It says that systems with feedback can progressively construct their own future possibility spaces, so that their eventual organization can contain information, differentiation, and apparent purpose that was nowhere represented at their origin.

The grades attach differently to the schema and to its instances. The feedback structure is established in the gravitational case, established as a framework in the geological case, and established as a phenomenon in the biological case. In the chemical case it is a plausible mechanism, and in the cultural case an analogy. The general schema, that some histories modify the mechanisms that determine what becomes possible next, is the essay's synthesis, supported by these instances and not proved by them.

\section{What follows}

If the argument holds, several consequences follow, each stated at the strength the preceding sections support.
\begin{itemize}
\item For origin-of-life research, emphasis shifts from the appearance of a first replicator to the appearance of closed positive causal paths and to how environments altered by such paths changed later chemistry. This is a suggestion about emphasis and does not decide between metabolism-first and replicator-first accounts.
\item For cosmology, the structure of a galaxy neither supports nor excludes an inference to a designer by itself. Any such inference needs independent grounds.
\item For biblical studies, literary parallels and historical existence are separate questions, and the productive question for each episode is which evidence is independent of the literary system.
\item For moral philosophy, the present justification of a value is not evidence of its origin, and the origin has to be reconstructed from history.
\end{itemize}

\section{Limits and open questions}

Several claims above are deliberately modest, and the following limits should be kept in view.
\begin{itemize}
\item The literary readings of Job, Mark, and the garment motif are statements about narrative mechanics. They do not settle historical questions, and the separation of literary dependence from historicity is meant to prevent them from being read that way.
\item Carrier's numerical conclusions are contested, and the argument does not rely on them.
\item The tidal-cycling hypothesis is contested and is not load-bearing.
\item Equation~\eqref{eq:parallel} is valid only for approximately independent trials, and the essay itself argues that independence fails once self-facilitating networks appear. No estimate of $N$ for the early Earth is offered.
\item The environment-mediated form of autocatalytic feedback is an extension of the formal theory and has not been shown here to occur in prebiotic chemistry.
\item The extension from chemistry to culture is an analogy of structure, and cultural transmission lacks a fixed replicator.
\end{itemize}

What the essay claims to have established is limited to three things. The inference from organized outcome to prior plan is not licensed by the outcome alone. Systems with outcome-to-condition feedback differ in kind from sampling processes, and an explanation of them must track how each stage altered the next. And the portable question of the first section can be asked of any case and has an answer that is determined by evidence, which in some cases includes a planner.

\appendix
\section{Evidence status of principal claims}

\begin{longtable}{p{0.44\textwidth}p{0.17\textwidth}p{0.29\textwidth}}
\hline
Claim & Grade & Basis \\
\hline
\endhead
Job 1 repeats a lone-survivor formula four times & Established & Job 1:15--19 \\
The earliest complete manuscripts of Mark end at 16:8 & Established & \cite{metzger} \\
Matthew and Luke narrate women who transmit, and Matthew narrates Galilee & Established & Matt.\ 28; Luke 24 \\
Later endings were composed to repair a felt gap & Plausible mechanism & Inference from the textual record \\
Mark 14:51--52 and Genesis 39:12 share a seizure and retained-garment pattern & Established (parallel); illustrative analogy (dependence) & Texts; \cite{scroggs} \\
Identity is carried by exterior tokens in Mark and in \emph{The Goat} & Illustrative analogy & \cite{goat} \\
Uncertainty about the crucifixion became an interpretive option in Muslim thought & Established & \cite{lawson} \\
Literary construction and historical existence are separate variables & Synthesis (method) & Reasoning from cases \\
Carrier's reference-class construction and probabilities & Contested & \cite{carrier1,carrier2,gullotta} \\
Cultural mechanisms of retention and repair are candidate sources of narrative organization & Plausible mechanism & \cite{lord,boyd,mesoudi} \\
Gravity yields a web of concentrations and voids from small perturbations & Established & \cite{bond} \\
Voids are underdense and contain galaxies & Established & \cite{vdw} \\
Angular momentum arises through tidal torques and needs asymmetry & Established & \cite{peebles,white} \\
Symmetric collapse and rebound in simple experiments & Illustrative analogy & Uncited simple simulation \\
Mineral inventory evolves through staged history & Established as framework & \cite{hazen2008} \\
Minerals adsorb, concentrate, and scaffold prebiotic molecules & Established in the laboratory & \cite{hazen2010} \\
These processes suffice for the origin of life & Plausible mechanism & \cite{hazen2010} \\
Fast tidal cycling favored nucleic-acid replication & Plausible mechanism, contested & \cite{lathe,varga} \\
The membraneless planetary cell picture & Illustrative analogy & None \\
Shoreline abrasion supplies relevant fresh surface & Conjecture & None \\
Equation~\eqref{eq:parallel} & Established mathematics; application conditional & Standard probability \\
Autocatalytic sets are well defined and can arise in models & Established within models & \cite{kauffman,hordijk2004,hordijk2019} \\
Prebiotic chemistry realized autocatalytic sets & Plausible mechanism & \cite{hordijk2019} \\
Environment-mediated feedback as autocatalysis & Synthesis (extension) & This essay \\
Niche construction modifies later selection & Established phenomenon; status as a distinct process debated & \cite{oslf} \\
Chemistry, biology, and culture iterate one recursion & Illustrative analogy & \cite{sperber} for the disanalogy \\
Genealogy forbids projecting a finished value backward into its origin & Established as a reading of Nietzsche & \cite{nietzsche} \\
History changes its own space of admissible continuations & Synthesis & This essay \\
\hline
\end{longtable}

\begin{thebibliography}{99}
\bibitem{bond} Bond, J. R., Kofman, L., and Pogosyan, D. (1996). How filaments of galaxies are woven into the cosmic web. \emph{Nature} 380, 603--606.
\bibitem{boyd} Boyd, R., and Richerson, P. J. (1985). \emph{Culture and the Evolutionary Process}. University of Chicago Press.
\bibitem{carrier1} Carrier, R. (2012). \emph{Proving History: Bayes's Theorem and the Quest for the Historical Jesus}. Prometheus Books.
\bibitem{carrier2} Carrier, R. (2014). \emph{On the Historicity of Jesus: Why We Might Have Reason for Doubt}. Sheffield Phoenix Press.
\bibitem{goat} Keaton, B., and St.\ Clair, M. (directors) (1921). \emph{The Goat}. Comedy short.
\bibitem{gullotta} Gullotta, D. N. (2017). On Richard Carrier's doubts: a response to Richard Carrier's \emph{On the Historicity of Jesus}. \emph{Journal for the Study of the Historical Jesus} 15(2--3), 310--346.
\bibitem{hazen2008} Hazen, R. M., Papineau, D., Bleeker, W., Downs, R. T., Ferry, J. M., McCoy, T. J., Sverjensky, D. A., and Yang, H. (2008). Mineral evolution. \emph{American Mineralogist} 93, 1693--1720.
\bibitem{hazen2013} Hazen, R. M. (2013). Paleomineralogy of the Hadean Eon: a preliminary species list. \emph{American Journal of Science} 313, 807--843.
\bibitem{hazen2010} Hazen, R. M., and Sverjensky, D. A. (2010). Mineral surfaces, geochemical complexities, and the origins of life. \emph{Cold Spring Harbor Perspectives in Biology} 2, a002162.
\bibitem{head} Head, P. M. A case against the longer ending of Mark. \emph{Text and Canon Institute}, \url{https://textandcanon.org/a-case-against-the-longer-ending-of-mark/}.
\bibitem{hordijk2004} Hordijk, W., and Steel, M. (2004). Detecting autocatalytic, self-sustaining sets in chemical reaction systems. \emph{Journal of Theoretical Biology} 227, 451--461.
\bibitem{hordijk2019} Hordijk, W. (2019). A history of autocatalytic sets. \emph{Biological Theory} 14(4), 224--246.
\bibitem{kauffman} Kauffman, S. A. (1986). Autocatalytic sets of proteins. \emph{Journal of Theoretical Biology} 119, 1--24.
\bibitem{lathe} Lathe, R. (2004). Fast tidal cycling and the origin of life. \emph{Icarus} 168, 18--22.
\bibitem{lawson} Lawson, T. (2009). \emph{The Crucifixion and the Qur'an: A Study in the History of Muslim Thought}. Oneworld.
\bibitem{lord} Lord, A. B. (1960). \emph{The Singer of Tales}. Harvard University Press.
\bibitem{m2m} Flyxion (2026). From minerals to minds: irreversibility and the physical origins of intelligence. \url{https://standardgalactic.github.io/alphabet/From\%20Minerals\%20to\%20Minds.pdf}.
\bibitem{mesoudi} Mesoudi, A. (2011). \emph{Cultural Evolution}. University of Chicago Press.
\bibitem{metzger} Metzger, B. M. (1994). \emph{A Textual Commentary on the Greek New Testament}, 2nd ed. United Bible Societies.
\bibitem{nietzsche} Nietzsche, F. (1998). \emph{On the Genealogy of Morality}, trans.\ M. Clark and A. J. Swensen. Hackett.
\bibitem{oslf} Odling-Smee, F. J., Laland, K. N., and Feldman, M. W. (2003). \emph{Niche Construction: The Neglected Process in Evolution}. Princeton University Press.
\bibitem{peebles} Peebles, P. J. E. (1969). Origin of the angular momentum of galaxies. \emph{Astrophysical Journal} 155, 393.
\bibitem{scroggs} Scroggs, R., and Groff, K. I. (1973). Baptism in Mark: dying and rising with Christ. \emph{Journal of Biblical Literature} 92, 531--548.
\bibitem{snapp} Snapp, J. E. A case for the longer ending of Mark. \emph{Text and Canon Institute}, \url{https://textandcanon.org/a-case-for-the-longer-ending-of-mark/}.
\bibitem{sperber} Sperber, D. (1996). \emph{Explaining Culture: A Naturalistic Approach}. Blackwell.
\bibitem{vdw} van de Weygaert, R., and Platen, E. (2011). Cosmic voids: structure, dynamics and galaxies. \emph{International Journal of Modern Physics: Conference Series} 1, arXiv:0912.2997.
\bibitem{varga} Varga, P., et al. (2006). Comment on the paper ``Fast tidal cycling and the origin of life'' by Richard Lathe. \emph{Icarus} 180, 274.
\bibitem{white} White, S. D. M. (1984). Angular momentum growth in protogalaxies. \emph{Astrophysical Journal} 286, 38.
\end{thebibliography}

\end{document}

